% GENERATED FILE. Edit canonical lessons, then run npm run generate:books.
% canonical-source-hash: fnv1a64:815b8ceafa99b21f
% canonical-lessons: MR-C355-matdan, MR-C355-pratinidhi, MR-C355-umedvar, MR-C355-pracar, MR-C355-nivdun-yene

\chapter{Election Day}
\label{ch:mr-election-day}
\hlchaptermodality{355}

\begin{chapteropening}
\textbf{By the end of this chapter:} \emph{I can talk about voting, candidates and representatives.}

Complete the last lesson of chapter 355: I can talk about voting, candidates and representatives.
\end{chapteropening}

\section[\texorpdfstring{\mr{मतदान} (matdān)}{मतदान (matdān)}]{\mr{मतदान} (matdān) — voting; a poll}

\label{lesson:MR-C355-matdan}



\begin{quote}
Before the new one: say the Marathi for population, then the Marathi for rural.
\end{quote}

\subsection*{You'll want to know: \mr{मतदान}}
\textbf{\mr{मतदान}} — \emph{matdān} — \textquotedblleft{}voting; a poll\textquotedblright{}.

\begin{quote}
\mr{रविवारी} \mr{सकाळपासून} \mr{मतदान} \mr{सुरू} \mr{झालं}.
\end{quote}

\begin{quote}
\emph{ravivārī sakāḷpāsūn matdān surū jhālaṁ.}
\end{quote}

\begin{quote}
Voting began on Sunday morning.
\end{quote}

Neuter.

\mr{मतदान} \mr{करणे}, to vote; \mr{मतदान} \mr{केंद्र}, a polling station.

\begin{cousinweb}
\textbf{\mr{मतदान}} (\emph{matdān}) is a learned Marathi compound built from Sanskrit \emph{mata}, \textquotedblleft{}opinion\textquotedblright{}, and \emph{dāna}, \textquotedblleft{}giving\textquotedblright{} (Monier-Williams): the giving of one's opinion.
\end{cousinweb}

\begin{grammarlens}[title={What you've built}]
One more word for giving reasons and opinions.
\end{grammarlens}

\subsection*{Your turn}
\begin{itemize}
\raggedright
  \item \emph{Say it:} \emph{matdān}
  \item \emph{Say it:} \emph{matdān}, once more
  \item \emph{Say it:} \emph{ravivārī sakāḷpāsūn matdān surū jhālaṁ.} — Voting began on Sunday morning.
  \item \emph{Recall:} say the Marathi for population, then the Marathi for rural, then say \emph{matdān} again
\end{itemize}

\begin{tcolorbox}[breakable,colback=teal!4,colframe=teal!35!black,title={Before you move on}]
What does \mr{मतदान} mean? (\textquotedblleft{}Voting; a poll\textquotedblright{}.) Say it once more.
\end{tcolorbox}
\section[\texorpdfstring{\mr{प्रतिनिधी} (pratinidhī)}{प्रतिनिधी (pratinidhī)}]{\mr{प्रतिनिधी} (pratinidhī) — a representative; a delegate}

\label{lesson:MR-C355-pratinidhi}



\begin{quote}
Before the new one: say the Marathi for rural, then the Marathi for voting.
\end{quote}

\subsection*{You'll want to know: \mr{प्रतिनिधी}}
\textbf{\mr{प्रतिनिधी}} — \emph{pratinidhī} — \textquotedblleft{}a representative; a delegate\textquotedblright{}.

\begin{quote}
\mr{आमच्या} \mr{भागाचा} \mr{प्रतिनिधी} \mr{लोकांचे} \mr{प्रश्न} \mr{मांडतो}.
\end{quote}

\begin{quote}
\emph{āmacyā bhāgācā pratinidhī lokāñce praśna māṇḍto.}
\end{quote}

\begin{quote}
Our area's representative puts forward people's problems.
\end{quote}

Masculine.

Used for a woman too: \mr{ती} \mr{आमची} \mr{प्रतिनिधी} \mr{आहे}; \mr{लोकप्रतिनिधी}, an elected representative.

\begin{cousinweb}
\textbf{\mr{प्रतिनिधी}} (\emph{pratinidhī}) is a tatsama, taken from Sanskrit \emph{pratinidhi}, \textquotedblleft{}a substitute, a representative\textquotedblright{} (Monier-Williams). Marathi lengthens the final vowel.
\end{cousinweb}

\begin{grammarlens}[title={What you've built}]
One more word for giving reasons and opinions.
\end{grammarlens}

\subsection*{Your turn}
\begin{itemize}
\raggedright
  \item \emph{Say it:} \emph{pratinidhī}
  \item \emph{Say it:} \emph{pratinidhī}, once more
  \item \emph{Say it:} \emph{āmacyā bhāgācā pratinidhī lokāñce praśna māṇḍto.} — Our area's representative puts forward people's problems.
  \item \emph{Recall:} say the Marathi for rural, then the Marathi for voting, then say \emph{pratinidhī} again
\end{itemize}

\begin{tcolorbox}[breakable,colback=teal!4,colframe=teal!35!black,title={Before you move on}]
What does \mr{प्रतिनिधी} mean? (\textquotedblleft{}A representative; a delegate\textquotedblright{}.) Say it once more.
\end{tcolorbox}
\section[\texorpdfstring{\mr{उमेदवार} (umedvār)}{उमेदवार (umedvār)}]{\mr{उमेदवार} (umedvār) — a candidate}

\label{lesson:MR-C355-umedvar}



\begin{quote}
Before the new one: say the Marathi for voting, then the Marathi for a representative.
\end{quote}

\subsection*{You'll want to know: \mr{उमेदवार}}
\textbf{\mr{उमेदवार}} — \emph{umedvār} — \textquotedblleft{}a candidate\textquotedblright{}.

\begin{quote}
\mr{या} \mr{जागेसाठी} \mr{चार} \mr{उमेदवार} \mr{उभे} \mr{आहेत}.
\end{quote}

\begin{quote}
\emph{yā jāgesāṭhī cār umedvār ubhe āhet.}
\end{quote}

\begin{quote}
Four candidates are standing for this seat.
\end{quote}

Masculine.

\mr{उभा} \mr{राहणे}, to stand for election; also a job applicant.

\begin{cousinweb}
\textbf{\mr{उमेदवार}} (\emph{umedvār}) is a loan from Persian \emph{ummedvār}, \textquotedblleft{}hopeful, a candidate\textquotedblright{}, built on \emph{ummed}, \textquotedblleft{}hope\textquotedblright{}.
\end{cousinweb}

\begin{grammarlens}[title={What you've built}]
One more word for giving reasons and opinions.
\end{grammarlens}

\subsection*{Your turn}
\begin{itemize}
\raggedright
  \item \emph{Say it:} \emph{umedvār}
  \item \emph{Say it:} \emph{umedvār}, once more
  \item \emph{Say it:} \emph{yā jāgesāṭhī cār umedvār ubhe āhet.} — Four candidates are standing for this seat.
  \item \emph{Recall:} say the Marathi for voting, then the Marathi for a representative, then say \emph{umedvār} again
\end{itemize}

\begin{tcolorbox}[breakable,colback=teal!4,colframe=teal!35!black,title={Before you move on}]
What does \mr{उमेदवार} mean? (\textquotedblleft{}A candidate\textquotedblright{}.) Say it once more.
\end{tcolorbox}
\section[\texorpdfstring{\mr{प्रचार} (pracār)}{प्रचार (pracār)}]{\mr{प्रचार} (pracār) — canvassing; publicity; propaganda}

\label{lesson:MR-C355-pracar}



\begin{quote}
Before the new one: say the Marathi for a representative, then the Marathi for a candidate.
\end{quote}

\subsection*{You'll want to know: \mr{प्रचार}}
\textbf{\mr{प्रचार}} — \emph{pracār} — \textquotedblleft{}canvassing; publicity; propaganda\textquotedblright{}.

\begin{quote}
\mr{उमेदवारांनी} \mr{घरोघरी} \mr{जाऊन} \mr{प्रचार} \mr{केला}.
\end{quote}

\begin{quote}
\emph{umedvārānnī gharoghrī jāūn pracār kelā.}
\end{quote}

\begin{quote}
The candidates canvassed from door to door.
\end{quote}

Masculine.

\mr{प्रचार} \mr{करणे}, to campaign; also to spread an idea.

\begin{cousinweb}
\textbf{\mr{प्रचार}} (\emph{pracār}) is a tatsama, taken from Sanskrit \emph{pracāra}, \textquotedblleft{}going about; coming forth, currency\textquotedblright{} (Monier-Williams).
\end{cousinweb}

\begin{grammarlens}[title={What you've built}]
One more word for giving reasons and opinions.
\end{grammarlens}

\subsection*{Your turn}
\begin{itemize}
\raggedright
  \item \emph{Say it:} \emph{pracār}
  \item \emph{Say it:} \emph{pracār}, once more
  \item \emph{Say it:} \emph{umedvārānnī gharoghrī jāūn pracār kelā.} — The candidates canvassed from door to door.
  \item \emph{Recall:} say the Marathi for a representative, then the Marathi for a candidate, then say \emph{pracār} again
\end{itemize}

\begin{tcolorbox}[breakable,colback=teal!4,colframe=teal!35!black,title={Before you move on}]
What does \mr{प्रचार} mean? (\textquotedblleft{}Canvassing; publicity; propaganda\textquotedblright{}.) Say it once more.
\end{tcolorbox}
\section[\texorpdfstring{\mr{निवडून} \mr{येणे} (nivḍūn yeṇe)}{निवडून येणे (nivḍūn yeṇe)}]{\mr{निवडून} \mr{येणे} (nivḍūn yeṇe) — to be elected}

\label{lesson:MR-C355-nivdun-yene}



\begin{quote}
Before the new one: say the Marathi for a candidate, then the Marathi for canvassing.
\end{quote}

\subsection*{You'll want to know: \mr{निवडून} \mr{येणे}}
\textbf{\mr{निवडून} \mr{येणे}} — \emph{nivḍūn yeṇe} — \textquotedblleft{}to be elected\textquotedblright{}.

\begin{quote}
\mr{आमच्या} \mr{गावचा} \mr{तरुण} \mr{निवडून} \mr{आला}.
\end{quote}

\begin{quote}
\emph{āmacyā gāvcā taruṇ nivḍūn ālā.}
\end{quote}

\begin{quote}
A young man from our village was elected.
\end{quote}

Past (he …): \mr{तो} \mr{निवडून} \mr{आला}, \emph{to nivḍūn ālā}.

Intransitive: the one elected is the subject; \mr{निवडून}, chosen, + \mr{येणे}, to come.

\begin{cousinweb}
\textbf{\mr{निवडून} \mr{येणे}} (\emph{nivḍūn yeṇe}) joins \mr{निवडून}, the -\mr{ऊन} form of \mr{निवडणे}, \textquotedblleft{}to choose\textquotedblright{}, an inherited verb, to \mr{येणे}, \textquotedblleft{}to come\textquotedblright{}; CDIAL lists Marathi \emph{nivaḍ}, \textquotedblleft{}choice, picked\textquotedblright{}, under \emph{nirvarayati} (CDIAL 7393).
\end{cousinweb}

\begin{grammarlens}[title={What you've built}]
One more word for giving reasons and opinions.
\end{grammarlens}

\subsection*{Your turn}
\begin{itemize}
\raggedright
  \item \emph{Say it:} \emph{nivḍūn yeṇe}
  \item \emph{Say it:} \emph{nivḍūn yeṇe}, once more
  \item \emph{Say it:} \emph{āmacyā gāvcā taruṇ nivḍūn ālā.} — A young man from our village was elected.
  \item \emph{Recall:} say the Marathi for a candidate, then the Marathi for canvassing, then say \emph{nivḍūn yeṇe} again
\end{itemize}

\begin{tcolorbox}[breakable,colback=teal!4,colframe=teal!35!black,title={Before you move on}]
What does \mr{निवडून} \mr{येणे} mean? (\textquotedblleft{}To be elected\textquotedblright{}.) Say it once more.
\end{tcolorbox}
